%%%%%%%%%%%%%%%%%%%%%%%%%%%%%%%%%%%%%%%%%%%%%%%%%%%%%%%%%%%%%%%%%%%%%%
% Lớp 10 | Mệnh đề. Tập hợp | Mệnh đề
% Dạng 0D1?1-6: Áp dụng mệnh đề vào suy luận có lí
% Số câu: 4  (sinh tự động bởi tach_cau.py; từ nay thầy sửa trực tiếp file này)
%%%%%%%%%%%%%%%%%%%%%%%%%%%%%%%%%%%%%%%%%%%%%%%%%%%%%%%%%%%%%%%%%%%%%%

%%%%%%%%%%%%%%%% TRẮC NGHIỆM 4 PHƯƠNG ÁN %%%%%%%%%%%%%%%%

% nguồn: MaHoaNH_TracNghiem_Lop10
\begin{ex}%[0D1V1-6]%[ma:cafb7995]
	Một nhóm học sinh M, N, P, Q, R xếp thành một hàng dọc trước một quầy nước giải khát. Dưới đây là  các thông tin ghi nhận được từ các học sinh trên
	\begin{multicols}{2}
		\begin{itemize}
			\item M, P, R là nam; N, Q  là nữ;
			\item M đứng  trước Q;
			\item N đứng  ở vị trí thứ nhất hoặc thứ hai;
			\item Học sinh đứng sau cùng  là nam.
		\end{itemize}
	\end{multicols}
	Hai vị  trí nào sau đây phải là hai học sinh khác giới tính?
	\choice 
		{Thứ hai và ba} 
		{Thứ hai  và năm}
		{\True Thứ ba và  tư}
		{Thứ ba và  năm}
	 \loigiai{
		Xét hai  trường hợp:
		\begin{itemize}
			\item Trường hợp 1. N đứng  ở vị trí thứ nhất.\\
			 Do M phải đứng trước Q nên M không thể xếp cuối cùng do đó có các trường hợp xếp:
			\begin{itemize}
				\item N đứng đầu, M, Q đứng tiếp theo và hai bạn sau  cùng là P, R
				\begin{center}
					\begin{tabular}{|c|c|c|c|c|}
						 \hline
						N & M & Q & P &  R \\
						 \hline
					\end{tabular} hoặc \begin{tabular}{|c|c|c|c|c|}
					\hline 
					N & M & Q & R & P \\ 
					 \hline
				\end{tabular}
			\end{center}
			Ta loại đáp án \lq\lq Thứ hai và  năm\rq\rq.
			\item N đứng đầu, xếp tiếp theo là một bạn nam (P hoặc R), M, Q, xếp cuối là một bạn nam còn lại 
			\begin{center}
				\begin{tabular}{|c|c|c|c|c|}
					\hline 
					N & P & M & Q &  R \\
					\hline 
				\end{tabular} hoặc \begin{tabular}{|c|c|c|c|c|}
				\hline 
				N & R & M &  Q & P \\
				\hline 
			\end{tabular}
		\end{center}
		Ta loại đáp án \lq\lq Thứ ba và năm\rq\rq. 
	\end{itemize}
	\item Trường hợp 2. N  đứng ở vị trí thứ hai.\\
	Do M phải đứng  trước Q nên M, Q phải xếp giữa N và học sinh nam cuối hàng.
	\begin{center}
		\begin{tabular}{|c|c|c|c|c|}
			\hline 
			R & N & M & Q & P \\ 
			\hline 
		\end{tabular} hoặc \begin{tabular}{|c|c|c|c|c|}
		\hline 
		P & N & M & Q  & R \\
		\hline 
	\end{tabular}
	\end{center}
	 Ta loại đáp án \lq\lq Thứ hai và ba\rq\rq.\\
	Vậy hai vị trí thứ ba và tư phải là hai  học sinh khác giới tính.
	\end{itemize}
 }
\end{ex}

% nguồn: MaHoaNH_TracNghiem_Lop10
\begin{ex}%[0D1C1-6]%[ma:54cf13e9]
	Một tòa  nhà có $n$ tầng, các tầng được đánh số từ $1$ đến $n$ theo thứ tự từ dưới lên. Có $4$ thang máy đang ở tầng $1$. Biết rằng mỗi thang máy có thể dừng ở đúng $3$ tầng và $3$ tầng này không là $3$ số nguyên liên tiếp và với hai tầng bất kỳ của tòa nhà luôn có một thang máy dừng được ở cả hai tầng này. Hỏi giá trị lớn nhất của $n$ là bao nhiêu?
	 \choice
		{\True  $6$}
		 {$7$}
		{$8$} 
		{$9$} 
	\loigiai{ 
		Giả  sử $4$ thang máy đó là $A,B,C,D$.\\
		Do khi bốc hai  thang bất kỳ luôn có một thang máy dừng được nên
		\begin{itemize}
			\item Khi bốc hai tầng $2,3$ có một thang dừng được giả sử đó là thang $A$, nên tầng $4$  không phải thang $A$ dừng.
			\item Khi bốc hai tầng $3,4$  có một thang dừng được giả sử đó là thang $B$, nên tầng $5$ không phải thang $B$ dừng.
			\item Khi bốc hai tầng $4,5$ có một thang dừng được giả sử đó là thang $C$, nên tầng  $6$ không phải thang $C$ dừng.
			\item Khi bốc hai tầng $5,6$ có một thang dừng được giả sử đó là thang $D$. 
			\item Khi bốc hai tầng $6,7$ có một thang dừng được khi đó không thể là thang $A,B,C$ vì  sẽ dừng $4$, thang $D$ không thể ở tầng $7$ do không thể ở ba tầng liên tiếp.
		\end{itemize}
		 Vậy khách sạn có tối đa sáu tầng.
	} 
\end{ex}


%%%%%%%%%%%%%%%% TRẢ LỜI NGẮN %%%%%%%%%%%%%%%%

% nguồn: MaHoaNH_TraLoiNgan_Lop10
\begin{bt}%[0D1H1-6]%[ma:deff5465]
	Cho phát biểu $\forall x\in\mathbb{R}$, $x^{2}+a-5>0$ với $a$ là số thực cho trước. Có bao giá trị  nguyên của $a<20$ để phát biểu đã cho là mệnh đề đúng?
	\shortans{14} 
	 \loigiai{
		Mệnh đề $\forall x\in\mathbb{R}$, $x^{2}+a-5>0$ đúng $\Leftrightarrow  x^{2}>5-a$ đúng $\forall x\in\mathbb{R}$.\\
		Điều này xảy ra khi và chỉ khi  $\max(5-a, 0) < 0$.\\
		Vì $x^2 \ge 0$, nên ta cần  $5-a$ phải nhỏ hơn $0$.\\
		 $5-a<0\Leftrightarrow a>5$.\\
		Kết hợp  với điều kiện $a<20$, ta có $5<a<20$.\\
		Các giá trị nguyên của $a$ là $6, 7, 8, \dots,  19$.\\
		Số giá trị nguyên của  $a$ là $19-6+1=14$.
	} 
\end{bt}


%%%%%%%%%%%%%%%% TỰ LUẬN %%%%%%%%%%%%%%%%

% nguồn: MaHoaNH_TuLuan_Lop10
\begin{bt}%[0D1V1-6]%[ma:98119397]
	Biết rằng có đúng $ 1  $ trong $ 4 $ mệnh đề sau là sai
	\begin{enumerate}[$ (1) $]
		\item \lq\lq  $A$ kém $B$ là $ 2 $ tuổi và $A$ hơn $C$ là $ 1 $ tuổi\rq\rq.
		\item \lq\lq $B$ không quá $ 25 $ tuổi và $B, C$ cách nhau $ 3 $  tuổi\rq\rq.
		 \item \lq\lq $C$ nhỏ tuổi hơn $A$ và $B$ hơn $A$ là $ 3 $ tuổi\rq\rq.
		\item \lq\lq $A$ là $ 22 $ tuổi hoặc $C$  là $ 25 $ tuổi\rq\rq.
	\end{enumerate}
	Hãy tính tổng  số tuổi của $A$, $B$ và $C$.
	\loigiai{ 
		Do theo $ (1) $ thì $A$ kém hơn $B$ là $ 2 $ tuổi, còn theo $ (3) $ thì $B$ hơn $A$  là $ 3 $ tuổi nên $ (1) $ và $ (3) $ không thể nào cùng đúng.\\
		Theo $ (2) $ thì $B$, $C$ cách  nhau $ 3 $ tuổi, còn theo $ (3) $ thì $B$ hơn $A$ là 3 tuổi nên $C$ hơn hoặc bằng tuổi $A$. Trái với $C$ nhỏ tuổi hơn $A$.\\
		 Vậy $ (3) $ sai, còn $ (1) $, $ (2) $ và $ (4) $ đúng.
		\begin{itemize}
			\item Nếu $A$  là $ 22 $ tuổi thì theo $ (1) $, ta có $B$ là $ 24 $ tuổi và $C$ là $ 21 $ tuổi. Thỏa mãn $ (2) $.
			\item Nếu $C$ là $ 25 $ tuổi thì theo $ (1) $ ta  có $A$ là $ 26 $ tuổi và $B$ là $ 28 $ tuổi. Không thỏa mãn $ (2) $.
		\end{itemize}
		Vậy $A$ là $22$ tuổi, $B$ là $24$ tuổi, $C$ là $21$ tuổi nên tổng số tuổi của $A$, $B$, $C$ là  $ 67 $ tuổi.
	} 
\end{bt}
